%=============================================================
% Chapitre 2 — Données et indicateurs
%=============================================================
\chapter{Données et indicateurs}
\label{ch:donnees}

\section{Sources de données}

Les données proviennent exclusivement des fichiers de recensement publiés par l'INSEE :

\begin{table}[H]
\centering
\caption{Sources de données utilisées dans le projet}
\label{tab:sources}
\begin{tabular}{llll}
\toprule
\textbf{Source} & \textbf{Année} & \textbf{Thème} & \textbf{Granularité} \\
\midrule
INSEE Filosofi / RP & 2016 & Population, emploi, logement & Canton \\
INSEE Filosofi / RP & 2022 & Population, emploi, logement & Canton \\
INSEE Filosofi & 2022 & Ménages, revenus & Canton \\
Ministère Intérieur & 2022 & Résultats électoraux & Département \\
\bottomrule
\end{tabular}
\end{table}

\section{Structure du jeu de données}

Après le processus d'ETL et d'ingénierie des variables, le jeu de données final
comprend :

\begin{itemize}
  \item \textbf{40 000 lignes} (cantons $\times$ tirages aléatoires pour équilibrer les classes)
  \item \textbf{171 colonnes} : variables brutes 2016, variables brutes 2022,
    variables delta, variable cible
  \item \textbf{27 variables delta} ($\Delta$-features) issues du calcul inter-recensements
  \item \textbf{5 classes} pour la variable cible, distribuées de façon synthétique
\end{itemize}

\section{Ingénierie des \texorpdfstring{$\Delta$}{Delta}-variables}

La clé de l'approche réside dans le calcul de variables d'évolution entre les deux
recensements (2016 et 2022). Pour chaque indicateur $X$, la variable delta est définie
par :

\begin{equation}
\Delta_X = \frac{X_{2022} - X_{2016}}{|X_{2016}| + \varepsilon}
\label{eq:delta}
\end{equation}

avec $\varepsilon = 10^{-8}$ pour éviter les divisions par zéro. Cette formulation
normalise automatiquement l'amplitude du changement par rapport à la valeur initiale.

\subsection{Variables delta principales}

Les 27 variables delta couvrent plusieurs domaines :

\begin{table}[H]
\centering
\caption{Variables delta principales (extrait)}
\label{tab:delta}
\begin{tabular}{lll}
\toprule
\textbf{Variable delta} & \textbf{Signification} & \textbf{Période} \\
\midrule
\texttt{delta\_P16\_POP}    & Évolution population totale     & 2016--2022 \\
\texttt{delta\_P22\_POP}    & Population totale 2022          & 2022 \\
\texttt{delta\_P22\_POP1564}& Pop. active (15--64 ans)        & 2016--2022 \\
\texttt{delta\_P22\_ACT1564}& Actifs 15--64 ans               & 2016--2022 \\
\texttt{delta\_P22\_CHOM1564}& Chômeurs 15--64 ans             & 2016--2022 \\
\texttt{delta\_P22\_EMPLT}  & Emplois total                   & 2016--2022 \\
\texttt{delta\_P22\_MEN}    & Ménages                         & 2016--2022 \\
\texttt{delta\_P22\_LOG}    & Logements totaux                & 2016--2022 \\
\texttt{delta\_P22\_LOGVAC} & Logements vacants               & 2016--2022 \\
\texttt{delta\_P22\_POP0014}& Population 0--14 ans            & 2016--2022 \\
\multicolumn{3}{c}{\textit{[17 autres variables delta de même nature]}} \\
\bottomrule
\end{tabular}
\end{table}

\section{Distribution de la variable cible}

La variable cible est construite synthétiquement à partir d'un score composite
socio-économique, puis discrétisée en 5 classes selon les quantiles imposés :

\begin{table}[H]
\centering
\caption{Distribution de la variable cible dans le jeu de données final}
\label{tab:distribution_cible}
\begin{tabular}{lrr}
\toprule
\textbf{Classe} & \textbf{Effectif} & \textbf{Proportion} \\
\midrule
Crise      & 4 000 & 10,0\% \\
Déclin     & 8 000 & 20,0\% \\
Stable     & 16 000 & 40,0\% \\
Croissance & 8 000 & 20,0\% \\
Boom       & 4 000 & 10,0\% \\
\midrule
\textbf{Total} & \textbf{40 000} & \textbf{100\%} \\
\bottomrule
\end{tabular}
\end{table}

\begin{warnbox}[Note méthodologique]
La variable cible est \textbf{synthétique} : elle ne provient pas de données d'enquête
d'opinion mais est construite à partir d'un score composite basé sur les évolutions
d'emploi, de chômage, de population et de logements. Cette approche permet d'entraîner
les modèles ML sur un volume de données suffisant (40 000 observations) tout en
ancrant la variable dans une réalité économique mesurable.
\end{warnbox}

\section{Stockage et flux des données}

\textbf{MongoDB} constitue le socle de persistance unique de l'ensemble du projet.
Les données transitent à travers quatre collections successives, chacune correspondant
à une étape clé du pipeline :

\begin{table}[H]
\centering
\caption{Collections MongoDB du projet Electio-Analytics}
\label{tab:mongo_collections}
\begin{tabular}{lll}
\toprule
\textbf{Collection} & \textbf{Contenu} & \textbf{Alimentée par} \\
\midrule
\texttt{raw\_data}    & CSV bruts INSEE (2016 + 2022)          & Phase 1--2 (wget + PyMongo) \\
\texttt{data\_final}  & Features nettoyées + 27 delta-variables & Phase 3--5 (Pandas + PyMongo) \\
\texttt{predictions}  & Scores par candidat, département, modèle & Phase 7 (pipeline ML) \\
\texttt{models\_meta} & Hyperparamètres et métriques des modèles & Phase 7 (pipeline ML) \\
\bottomrule
\end{tabular}
\end{table}

\noindent Ce flux unidirectionnel garantit la traçabilité complète :
aucune étape aval ne lit des fichiers locaux, toutes les lectures se font
exclusivement via \texttt{PyMongo}.

%=============================================================
\section{Indicateurs absents du cahier des charges : justification}
\label{sec:indicateurs_absents}
%=============================================================

Le cahier des charges initial mentionnait quatre familles d'indicateurs
complémentaires : sécurité publique, vie associative, taux de pauvreté monétaire
et nombre d'entreprises. Leur absence dans la version livrée ne résulte pas d'un
oubli méthodologique, mais de contraintes objectives de granularité et de
disponibilité des données publiques.

\paragraph{Sécurité publique (SSMSI).}
Les statistiques de criminalité publiées par le Service Statistique Ministériel de
la Sécurité Intérieure sont disponibles par département ou par zone de ressort de
groupement de gendarmerie. La maille cantonale n'est pas proposée dans les fichiers
\texttt{data.gouv.fr} pour ces indicateurs. Intégrer des valeurs départementales
dans un modèle entraîné au canton aurait introduit une ambiguïté statistique
inacceptable : tous les cantons d'un même département auraient reçu la même valeur,
réduisant l'indicateur à un simple effet-département sans pouvoir discriminant.

\paragraph{Vie associative (RNA).}
Le Répertoire National des Associations est disponible à l'échelle de la commune,
non du canton. Une agrégation manuelle était techniquement envisageable mais aurait
nécessité un référentiel de correspondance commune-canton exhaustif et à jour ---
ce référentiel n'est pas fourni directement par l'INSEE et doit être reconstitué
à partir de la table de passage COG, opération hors périmètre du projet compte tenu
du calendrier imposé.

\paragraph{Taux de pauvreté monétaire (Filosofi).}
L'indicateur de pauvreté au seuil de 60\% du revenu médian est disponible dans
Filosofi à la maille IRIS et commune, mais \textbf{non systématiquement au canton}
pour les deux millésimes comparés (2016 et 2022). Les variables de ménages et
d'emplois déjà intégrées constituent un proxy partiel de la fragilité économique
sans nécessiter cette agrégation incertaine.

\paragraph{Nombre d'entreprises (SIRENE).}
Le répertoire SIRENE recense les établissements à la commune. L'agrégation au canton
est possible en théorie mais imprécise : certaines communes historiques chevauchent
plusieurs cantons dans les découpages antérieurs à 2015, créant des doublons ou
des omissions selon le référentiel de correspondance utilisé.

\begin{warnbox}[Perspectives d'enrichissement — version 2]
Une version enrichie pourrait intégrer : (1)~les données SSMSI à la maille
départementale comme variable contextuelle externe aux cantons ;
(2)~le nombre d'associations par canton agrégé depuis le RNA via l'API
\texttt{data.gouv.fr} et la table de passage COG ;
(3)~le taux de pauvreté Filosofi disponible à l'IRIS, agrégé par pondération
de population cantonale ;
(4)~les établissements actifs SIRENE convertis au canton via la table de passage
commune-canton publiée par l'INSEE depuis 2022.
\end{warnbox}

%=============================================================
\section{Données historiques de vote : justification du périmètre temporel}
\label{sec:vote_historique}
%=============================================================

Le cahier des charges demandait d'inclure les résultats électoraux des scrutins
présidentiels de 2012 et 2017 comme variables d'entrée ou de validation. Le jeu
de données final utilise uniquement les résultats de 2022 à des fins de validation
externe. Ce choix résulte d'une contrainte de cohérence géographique.

\paragraph{Problème d'agrégation commune-canton.}
Les résultats du premier tour des présidentielles 2012 et 2017 sont mis à
disposition par le Ministère de l'Intérieur à la maille \textbf{commune} sur
\texttt{data.gouv.fr}. Or notre modèle opère à la maille \textbf{canton}, définie
par le redécoupage territorial de 2015 qui regroupe en moyenne plusieurs communes.
Agréger les résultats communaux au canton aurait nécessité une pondération par
nombre d'inscrits sur les listes électorales --- opération techniquement réalisable
mais introduisant une incertitude supplémentaire dans des données qui servent à
\emph{valider} les prédictions, pas à les entraîner.

\paragraph{Choix de prioriser la qualité spatiale.}
Le projet vise à prédire le comportement électoral à partir d'indicateurs
socio-économiques, et \textbf{non} à construire un modèle autorégressif qui
prédirait 2022 à partir des votes 2017. Intégrer les résultats passés comme
features aurait risqué de créer une dépendance temporelle masquant l'apport réel
des indicateurs économiques --- ce qui aurait affaibli la démonstration centrale
du projet. Nous avons donc délibérément choisi d'entraîner le modèle uniquement
sur des indicateurs socio-économiques mesurables indépendamment du vote, et de
réserver les résultats 2022 à la validation finale par département.

\begin{infobox}[Perspective v2]
La table de passage commune-canton officielle publiée par l'INSEE (référentiel COG
2022) permettrait d'agréger les résultats 2017 avec une correspondance fiable.
L'ajout d'un indicateur de \og stabilité électorale \fg{} (part du vote 2017
reportée sur le même bloc politique en 2022) constituerait une feature pertinente
sans compromettre la cohérence spatiale du modèle.
\end{infobox}
